\section{Physics Layer III: Oscillation, Phase, and Synchronization}

The synchronization textbooks in the literature map matter because many
transitions are coupled to external timing structures. A schedule, a partner,
or an institutional clock can entrain behaviour.

\begin{discussion}
Synchronization is not primarily about social pressure. It is about phase
locking. If two or more processes repeatedly align, the target system loses
degrees of freedom to drift into arbitrary timing. That is why shared start
times, repeated class schedules, and co-working sessions often feel easier than
privately chosen starts: coupling reduces timing uncertainty.
\end{discussion}

\begin{example}[Co-working as phase control]
Two students agree to begin at 8:00 PM every evening. The arrangement is useful
not only because each fears disappointing the other. It is useful because the
timing of the state transition becomes externally coupled. The intervention
therefore acts on phase alignment, not only on motivation.
\end{example}
